% titel: Blatt 5, A7: Höhere 1. Ionisierungsenergie
% skript: Mortimer Kap. 7 (chemie: Ionisierungsenergie)

Die 1.\ Ionisierungsenergie ist die Energie, das erste Elektron eines
neutralen Atoms zu entfernen. Sie hängt von Kernladung, Schalenzahl und
Konfigurationsstabilität ab und folgt damit dem Periodensystem.

\begin{werkzeugkasten}[Ionisierungsenergie vergleichen]
\begin{itemize}
  \item \textbf{In der Periode} nach rechts grösser: steigende
        $Z_{\mathrm{eff}}$ zieht die Valenzelektronen fester an (Edelgas rechts
        mit der höchsten IE).
  \item \textbf{In der Gruppe} nach oben grösser: kleinere Atome, weniger
        Schalen $\to$ Valenzelektronen fester gebunden.
  \item Ausnahmen durch stabile halb-/vollbesetzte Unterschalen; die
        Edelgas-Dominanz (rechts/oben) bleibt aber der Leitfaden.
\end{itemize}
\textit{Skript:} Mortimer Kap. 7 (Periodizität).
\end{werkzeugkasten}

\begin{achtung}[Darauf achten]
\begin{itemize}
  \item IE-Trend ist \textit{gegenläufig} zum Radius-Trend: kleiner Radius
        $\leftrightarrow$ höhere IE. Beide Fragentypen nicht verwechseln.
  \item Entscheidend ist $Z_{\mathrm{eff}}$: wächst nach rechts, bleibt in der
        Gruppe annähernd gleich, während die Schale wächst.
  \item Bei Edelgas-Partnern (Se/Kr, Kr/Xe) spielen zwei Argumente zusammen
        (rechts bzw.\ oben) --- benennen, welches Paar-Hauptargument zählt.
\end{itemize}
\end{achtung}
